\scanpage{41}
\uline{Résultats de Kostant} sur classes de conjugaison remarquables dans $G$ $1/2$ simple adjoint (sur $\mathbb C$)

a) Classes de conjugaison de ss-groupes simples de rang $1$ $\deleted{\unsure{\simeq}}\unsure{\simeq}$ classes de conjugaison d'éléments unipotents de $G$ [ou de $\unsure{\mathfrak g|V}$\inserted{ou de ss-groupes additifs de $G$},] $\simeq$ certaines classes de conjugaison \deleted{d'éléments semi-simples de $\mathfrak g$ [\uline{NB} \unsure{en de hom}\illegible{1} \unsure{princ.} ss} d'homomorphismes $\mathbb G_{\mathrm m}\to G$ [ou d'éléments de $\mathfrak g$ \ldots].

Les hom admissibles $G_{\mathrm m}\to T\simeq\deleted{\unsure{T^{\Sigma}}}G_{\mathrm m}^{\Sigma}$ ($\Sigma$ ensemble des racines simples) dans le ``domaine fondamental'' $W$ de $\mathrm{Hom}(G_{\mathrm m},T)\simeq\mathbb Z^{\Sigma}$, savoir $\mathbb N^{\Sigma}$, correspondent à \inserted{des} éléments dont les composantes sont $0,1,2$ (pas tous)

% margin note written sideways
\margin{l'élément $w_0\in W$ qui transforme la chambre $C$ en l'opposée}

b) Le cas où toutes les composantes sont $1$ est admissible ; il correspond aux \uline{éléments unipotents principaux}\inserted{(ou \uline{réguliers})}, et dans $\mathfrak g$ aux \uline{éléments semi-simples principaux} ; \deleted{\unsure{aux}} \uline{hom} ``\uline{principaux}'' $\mathbb G_{\mathrm m}\to G$. On parle des sous-groupes simples de rang $1$ \uline{principaux}.

\deleted{\unsure{b)}}

Centralisateur d'un élément unipotent principal : sous-groupe \uline{connexe}\inserted{(\unsure{Springer|de Springer})} \inserted{\uline{unipotent}} \uline{commutatif}\inserted{\uline{de dim $\unsure{r}$}}

$\times$\quad Normalisateur d'un ss-groupe additif principal : groupe connexe \deleted{\unsure{de dim}}\inserted{\unsure{son quotient} par \unsure{précéd.} $\simeq G_{\mathrm m}$ (\unsure{canont})}

Centralisateur d'un hom principal $\mathbb G_{\mathrm m}\to G$ : un tore maximal. Donc ces hom sont en correspondance biunivoque avec les couples de Killing de $G$.

\margin{Donc sur base quelc $\exists$ ss-\deleted{\unsure{groupe simple}}\inserted{de} \deleted{\unsure{de}} rg $1$ princip. \uline{déployé} ssi $G$ quasi-déployé}

Normalisateur d'un $1$-tore principal : groupe $\deleted{\unsure{N}}N_0$ dans lequel $T$ est d'indice $2$, correspondant à \unsure{$b|w$}

\deleted{\unsure{Centralisateur}}\inserted{Normalisateur} d'un sous-groupe \deleted{principal} de rang $1$ principal : lui-même.

\deleted{Normalisateur}\inserted{Stabilisateur} d'un ss-gpe principal de rang $1$\inserted{$H$}, avec un tore maximal $T_0$ donné dans celui-ci (resp. un \deleted{couple de}\inserted{Borel, resp. un} Killing) : Normalisateur\inserted{$N_0$} du tore\inserted{$T_0$} dans $H$\unsure{.} (resp. $B_0$, resp. $T_0$)

\margin{\deleted{\unsure{Donc sur un corps} \illegible{1} couple \unsure{de rang} $1$}\inserted{\uline{NB} $\mathbb Z/2\mathbb Z$ opère sur \unsure{KillCl}, échangeant les couples opposés.} \deleted{\unsure{là-dessus, on a}\illegible{2}} \unsure{quotient est le schéma des} \deleted{\unsure{ss}} tores de dim $1$ \uline{principaux}}

c) Soit $G_{\mathrm m}\xrightarrow{u}G$ un hom principal, et soit \deleted{\illegible{1}} $h$\inserted{$=o(\psi)+1$} \deleted{\unsure{le}} le nb de Coxeter ($=$ somme des coeff. de la plus haute racine $\psi$, \uline{plus} $1$).

Soit $\zeta$ \deleted{\unsure{un él.|une}} racine primitive $h$.\textsuperscript{ème} de $1$, et considérons $u(\zeta)$. On prouve que les $u(\zeta)$ ($\zeta$ variable, et évidemment aussi pour $u$ variable) sont conjugués. Les éléments obtenus sont appelés \inserted{\uline{éléments}} \uline{principaux} \inserted{\uline{réguliers}} de $G$. Si $T$ est un tore maximal, et $\sigma\in W=N(T)/T$ \deleted{\unsure{élémt}} un \deleted{\unsure{élémt}} élément de Coxeter, on prouve que les représentants de $\sigma$ sont tous principaux (donc réguliers).

\margin{\unsure{en car qui|encore que|ordre qui} divise $h$, il vaut mieux considérer la classe d'hom. $\mu_h\to G$}

Centralisateur d'un élément régulier principal $\gamma$ : ss-gpe de $N(T)$ \deleted{\unsure{formé}} d'image réc.\ du ss-gpe de $W=W(T)$ qui stabilise $\gamma$.

\margin{\textcircled{?}} C'est probablement $T$, i.e. \deleted{\unsure{on a ainsi}} la donnée de $\gamma$ revient à celle de $\unsure{\varphi|u|\mu}:G_{\mathrm m}\to G$

\margin{$G$ \uline{simple}}
d) Considérons l'opération de $W$ sur $\deleted{\unsure{\mathfrak g}}$ $\underline t=\mathrm{Lie}(T)$. Il existe des \deleted{\unsure{dr. propres}}\inserted{droites} propres\inserted{$D$} \deleted{\unsure{pour|dans}} $\sigma$ (op. de Coxeter) dont la valeur propre est une racine primitive\inserted{$h$.\textsuperscript{ème}} de $1$, \inserted{ce sont exactement celles qui sont ``régulières'',} et ces droites sont conjuguées sous l'action de $W$ (donc : \unsure{fusion|faisceau|forment} conjuguées
% archive sheet number: 40

